%------------------------------
% DOCUMENT SETTINGS
%------------------------------
\documentclass[11pt]{exam}
\printanswers % Uncomment to show answers

%------------------------------
% PACKAGES
%------------------------------
\usepackage{amsmath,amsfonts,amssymb,amsthm} % math fonts and symbols
\usepackage{bbm} % math blackboard font
\usepackage{enumitem} % for customizable lists
\usepackage{textcomp} % symbols
\usepackage[top=2cm,bottom=2cm,right=1cm,left=1cm,paper=a4paper,headheight=6cm]{geometry} % margins
\usepackage{xcolor} % colors
\usepackage{graphicx} % figures
\usepackage{booktabs} % tables

%------------------------------
% MACROS AND COMMANDS
%------------------------------
\newcommand{\N}{\mathbb{N}} % natural numbers
\newcommand{\R}{\mathbb{R}} % real numbers

%------------------------------
% HEADER AND FOOTER
%------------------------------
\pagestyle{headandfoot}
% degree, subject, exam, author, date, university
\newcommand{\degree}{Bachelor in Philosophy, Politics and Economics}
\newcommand{\shortdeg}{PPE}
\newcommand{\exam}{Final Exam}
\newcommand{\subject}{Quantitative Methods for Humanities}
\newcommand{\theyear}{2024/25}
\newcommand{\ccauthor}{A.G. Azze}
\newcommand{\thedate}{May 22, 2025}
\newcommand{\university}{CUNEF Universidad}

\header{\degree \\ \subject{} -- \theyear}{}{\thedate \\ \ccauthor}
\footer{\university}{}{\thepage}

%------------------------------
% DOCUMENT
%------------------------------
\begin{document}

%--- Instructions ---
\begin{center}

    {\Huge \textbf{\exam{}}}
    \vspace{0.7cm}

    \begin{flushright}
        \textbf{Time Limit:} 120 minutes
    \end{flushright}
    \vspace{-0.2cm}

    \fbox{\parbox{\textwidth}{\centering 
        \begin{enumerate}[leftmargin = 0.5cm, label = $\bullet$, rightmargin = 0.2cm]
    
            \item \textbf{Every non-trivial calculation and claim in the exam must be properly justified.}

            \item Digital devices such as mobile phones, tablets, and laptops, as well as internet access, are forbidden.
            
        \end{enumerate}
    }}
    
\end{center}

%--- Case Study Scenario ---
Three social scientists conducted a randomized controlled trial to examine whether social pressure within neighborhoods increases voter participation\footnote{Alan S. Gerber, Donald P. Green, and Christopher W. Larimer (2008) “Social pressure and voter turnout: Evidence from a large-scale field experiment.” \emph{American Political Science Review}, vol. 102, no. 1, pp. 33–48.}. During the 2006 primary election in Michigan, registered voters were randomly assigned to one of four groups, each receiving a different get-out-the-vote postcard or no mailing:
\begin{itemize}[noitemsep]
  \item \emph{Civic Duty}: a standard civic-duty reminder emphasizing the importance of voting.
  \item \emph{Neighbors}: the same civic-duty text plus a notice that the recipients’ voting records would be reported to their neighbors.
  \item \emph{Hawthorne}: the same civic-duty text plus the opening ``\textbf{YOU ARE BEING STUDIED!}'' to highlight their participation on the study.
  \item \emph{No-Message}: no postcard sent.
\end{itemize}

Table \ref{tab:turnout} summarizes each group’s 2006 turnout rate and key pre-treatment variables:
\begin{table}[ht]
    \centering
    \caption{Voter Turnout per Group}
    \label{tab:turnout}
    \begin{tabular}{ccccc}
        \toprule
        \textbf{Group:} & Civic Duty & Neighbors & Hawthorne & No-Message \\
        \midrule
        \textbf{Turnout rate 2006:} & 0.2966 & 0.3145 & 0.3223 & 0.3779 \\
        \textbf{Average age:} & 49.8135 & 49.6590 & 49.7048 & 49.8529 \\
        \textbf{Turnout rate 2004:} & 0.4003 & 0.3994 & 0.4032 & 0.4066 \\
        \bottomrule
    \end{tabular}
\end{table}

\begin{questions}

% Question 1
\question[2] For a randomly selected voter in the 2006 Michigan primary election:
\begin{enumerate}[label=(\alph*)]
\item What is the probability that the voter:
  \begin{enumerate}[label=(\roman*)]
    \item did not receive any message;
    \item voted;
    \item both did not receive any message and voted.
    \item received no message provided that she voted.
  \end{enumerate}
\item Are the events defined in (i) and (ii)independent? Justify briefly.
\end{enumerate}

\begin{solution}
\begin{enumerate}[label=(\alph*)]
\item These are the required probabilities:
    \begin{enumerate}[label=(\roman*)]
        \item With four equally likely groups, $P(\text{No-Message})=1/4=0.25$.  
        \item Overall turnout is the mean of the group rates:
        \[
        P(\text{Voted})=\frac{0.3145+0.2966+0.3223+0.3779}{4} \approx 0.3278.
        \]
        \item $
        P(\text{No-Message}\cap\text{Voted}) 
        = P(\text{Voted} \mid \text{No-Message}) P(\text{No-Message}) 
        = 0.25\times0.3779 \approx 0.0945
        $.
        \item $
        P(\text{No-Message} \mid \text{Voted}) 
        = \frac{P(\text{Voted} \mid \text{No-Message})P(\text{No-Message})}{P(\text{Voted})} 
        = \frac{0.3779 / 4}{0.3278} \approx 0.2882
        $
    \end{enumerate}
    \item $P(\text{No-Message})P(\text{Voted}) = 0.25\times0.3278 \approx 0.0820 \neq 0.0945$, so they are not independent.
\end{enumerate}
\end{solution}

% Question 2
\question[2] To asses the effect of different social-pressure messages, answer the following questions.
\begin{enumerate}[label=(\alph*)]
    \item Specify the treatment variable and the outcome variable when estimating the effect of the Neighbors postcard.
    \item Using the 2006 turnout rates from Table \ref{tab:turnout}, compute the Difference-in-Means (DiM) estimator for Neighbors versus No-Message.
    \item Given the information of the average age and 2004 turnout rate in Table \ref{tab:turnout}, does randomization appear successful? What does this imply for the DiM estimator? Why?
    \item Compute DiM estimators for Civic Duty versus Hawthorne messages. What does this suggest about being observed?
\end{enumerate}

\begin{solution}
\begin{enumerate}[label=(\alph*)]
    \item 
    $T = 1 \text{ if individual got a Neighbors postcard}, = 0 \text{ if no message}$ \\
    $Y = 1 \text{ if the individual voted, } = 0 \text{ if abstained}$.
    % \item Yes. Each voter’s postcard assignment and response do not affect others’ behavior (postcards sent individually).
    \item 
    $DiM_{Neighbors} = 0.3145 - 0.3779 = -0.0634.$ Indicates a 6.34 percentage-point decrease in turnout under Neighbors vs. No-Message.
    %Since control (No-Message) rate is 0.3779, effect of social-pressure is 0.3145−0.3779=−0.0634 (i.e. a 6.34 pp decrease).
    \item The four groups have nearly identical age ($\approx 49.7$) and prior turnout ($\approx 0.40$), so randomization appears balanced, supporting unbiased DiM estimates due to equal distribution of potential confounders among the treatment and control groups.
    \item 
    $DiM_{Hawthorne-vs-CivicDuty} = 0.3223 - 0.2966 = 0.0257.$ Suggests that awareness of being studied increases turnout by about 2.6 percentage points over the civic-duty message.
    % \item The Hawthorne effect is $-0.0556$ (meaning −5.56 percentage points), which is smaller in magnitude than Civic Duty (−8.13 pp), suggesting that awareness of being studied partially offsets the civic reminder’s negative impact, but less so than social-duty framing.
\end{enumerate}
\end{solution}

% Question 3
\question[2] We now combine 2004 and 2006 turnout data to assess the Neighbors versus the No-Message effect:
\begin{enumerate}[label=(\alph*)]
  \item Identify the treatment and control groups, and the before and after periods
  \item Compute the Before-and-After (BA) estimator. Interpret.
  \item Under which assumption the BA estimator would be unbiased for the caussal effect you are trying to asses? Do you think it is reasonable to assume in this case?
  \item Obtain the Difference-in Differences (DiD) estimator.
  \item What is the main assumption behind the DiD estimator to guarantee its unbiasness.
\end{enumerate}

\begin{solution}
\begin{enumerate}[label=(\alph*)]
    \item Treatment group: Neighbors postcard; control group: No-Message. \\
          Before period: 2004 primary; after period: 2006 primary.
    
    \item $BA = \text{turnout}_{2006}^{\text{No-message}} - \text{turnout}_{2004}^{\text{No-message}} = 0.3145 - 0.3994 = -0.0849$. That is, there was a reduction of $8.49$ percentage points in turnout in the Neighbors group when comparing 2004 versus 2006 data.
    \item The BA estimator is unbiased if there are no confounders changing between 2004 and 2006. This is unlikely, since other factors (e.g., campaign intensity) may differ.
    \item $DiD = (0.3145 - 0.3994) - (0.3779 - 0.4066) = -0.0849 - (-0.0287) = -0.0562$. After adjusting the BA estimate by subtracting the trend change, the DiD estimate suggests a decrease in 5.62 percentage points in turnout rate.
    \item DiD validity requires the parallel trends assumption: in absence of treatment, Neighbors and No-Message groups would have experienced the same change in turnout.
\end{enumerate}
\end{solution}

% Question 2
\question[4] Given the estimated model
\[
\widehat{\texttt{turnout\_2006}} = 0.3145 - 0.0063\,\texttt{message\_neighbors} + 0.007\,\texttt{message\_hawthorne} - 0.018\,\texttt{message\_no-message},
\]
where \texttt{message\_neighbors}, \texttt{message\_hawthorne}, and \texttt{message\_no-message} are binary variables equal to 1 if the voter received that specific message and 0 otherwise. The fitted model has an $R^2$ of $0.0033$.
\begin{enumerate}[label=(\alph*)]
  \item Compute the predicted 2006 turnout rate for the civic-duty message group.
  \item Interpret the coefficient associated to \texttt{message\_hawthorne}.
  \item Can you regard the coefficient associated to \texttt{message\_no-message} as an average causal effect? Why or why not?
  \item Comment on the predictive power of the model.
\end{enumerate}
Suspecting that the Neighbors effect varies by prior turnout, you consider the interaction model
\[
\widehat{\texttt{turnout\_2006}} = \widehat{\beta}_0 + \widehat{\beta}_1\,\texttt{turnout\_2004} + \widehat{\beta}_2\,\texttt{message\_neighbors} + \widehat{\beta}_3\,(\texttt{turnout\_2004}\times\texttt{message\_neighbors}),
\]
with an $R^2$ of $0.3078$.
\begin{enumerate}[label=(\alph*)]
  \item[(e)] Use this model to quantify how much does the Neighbors effect change for someone who voted in 2004 compared to someone who did not.
  \item[(f)] If your sole goal was to predict the 2006 turnout rate, would you use this model or the previous one? Why?
\end{enumerate}

\begin{solution}
\begin{enumerate}[label=(\alph*)]
  \item Since all the binary variables included in the regression are zero for the civic-duty group, the estimated turnout rate for this group is the intercept, that is, $\widehat{\texttt{turnout\_2006}} = 0.3145$.
  \item The coefficient of the Hawthorne group ($0.007$) is the estimated increase in turnout of this group relative to Civic Duty.
  \item The coefficient of the No-message group can be regarded as the estimated causal effect of recieving a civic-duty message versus not receiving any message, since the study is a RCT
  \item The model's $R^2$ indicates that only about $0.3\%$ of the variability is being explained, which is a very poor performance for prediction purposes.
  \item For non-voters in 2004 (\texttt{turnout\_2004} $= 0$), the Neighbors effect is $\widehat{\beta}_2$. For prior voters (\texttt{turnout\_2004} $= 1$), it is $\widehat{\beta}_2+\widehat{\beta}_3$. Hence, the difference is $\widehat{\beta}_2$.
  \item The new interaction model is about 100 times more capable in terms of prediction, as its $R^2$ indicates that about $30\%$ of the variability is captured.
\end{enumerate}
\end{solution}

\end{questions}

\end{document}
